%%% FIELD proof-minimax-length-frontier
Put \(q=(1-\alpha)/2\).  Uniform coverage of \(I_{T,\Delta}\) over \(\mathfrak M_T(\alpha,\beta,L,c,C,\varepsilon)\) follows directly from \(\mathrm{thm:uniform\text{-}gated\text{-}inference}\), because that theorem is unconditional on quota success and is uniform over the union containing this class.  We prove the expected-length assertion and the two converse components.

\emph{Upper bound.}
By \(\mathrm{lem:quota\text{-}success}\), uniformly over the class,
\[
 K_1(n_T)\asymp T^{1-\alpha},\qquad
 K_0(n_T)\asymp T,
 \qquad J_1\asymp T^{\rho(1-\alpha)},\qquad
 J_0\asymp T^\rho. \tag{1}
\]
Thus, for the quantities in \(\mathrm{def:hybrid\text{-}interval}\),
\[
 S_T=K_1(n_T)^{-1}+K_0(n_T)^{-1}
     =O\!\left(T^{-(1-\alpha)}\right),
 \qquad
 a_T=\sqrt{2S_Tx_\eta}=O(T^{-q}). \tag{2}
\]
The cohort schedule gives \(g_{T,\Delta}=m_T+\Delta_T\geq m_T\asymp T^\rho\).  Hence the polynomial delay envelope in the model-class definition yields
\[
 v_{T,\Delta}
 =\min\{1,L(1+g_{T,\Delta})^{-\beta}\}
 =O(T^{-\rho\beta}). \tag{3}
\]
For every fixed \(x>0\), equations (1)--(3) and the definition
\(
A_{a,T,\Delta}(x)=\sqrt{2v_{T,\Delta}x/J_a}+4x/(3J_a)
\)
give
\[
 \begin{aligned}
 A_{1,T,\Delta}(x)
 &=O\!\left(T^{-\rho(\beta+1-\alpha)/2}
          +T^{-\rho(1-\alpha)}\right),\\
 A_{0,T,\Delta}(x)
 &=O\!\left(T^{-\rho(\beta+1)/2}+T^{-\rho}\right).
 \end{aligned} \tag{4}
\]
The declared restriction
\(
\rho>(1-\alpha)/(1-\alpha+\beta)
\)
implies \(\rho(\beta+1-\alpha)>1-\alpha\), while \(\rho>1/2\) implies \(\rho(1-\alpha)>(1-\alpha)/2=q\).  Also
\(
(1-\alpha)/(1-\alpha+\beta)\geq(1-\alpha)/(1+\beta)
\), so \(\rho(1+\beta)>1-\alpha\), and \(\rho>1/2\geq q\).  Consequently both lines of (4) are \(o(T^{-q})\).  These inequalities are strict, so the same calculation absorbs the factor \(\log T\) when \(x=4\log T\), proving
\[
 d_{T,\Delta}
 =A_{1,T,\Delta}(4\log T)+A_{0,T,\Delta}(4\log T)
 =o(T^{-q}). \tag{5}
\]

Every main-prefix variable lies in \([-1,1]\).  Since each quota is at least two, its usual unbiased sample variance is at most two.  Therefore, on \(E_T\),
\[
 \widehat s_{T,\Delta}^2\leq2S_T,
 \qquad \widehat s_{T,\Delta}=O(T^{-q}). \tag{6}
\]
Let \(z=\Phi^{-1}(1-\eta/2)\).  For \(u\geq0\), evaluating the defining coverage function of \(c_\eta(u)\) at \(u+z\) gives
\[
 \Phi(z)+\Phi(2u+z)-1\geq2\Phi(z)-1=1-\eta.
\]
That coverage function is increasing in its critical-value argument, so uniqueness in \(\mathrm{def:bounded\text{-}shift\text{-}critical\text{-}value}\) implies
\[
 c_\eta(u)\leq u+z,
 \qquad
 \widehat s_{T,\Delta}c_\eta(B_{T,\Delta}/\widehat s_{T,\Delta})
 \leq B_{T,\Delta}+z\widehat s_{T,\Delta} \tag{7}
\]
whenever \(\widehat s_{T,\Delta}>0\); the zero-standard-error case uses the concentration branch by definition.

Next, \(h_{T,\Delta}=H_T-k_T\), \(H_T\geq T\), and \(k_T=O(T^\rho)\) with \(\rho<1\).  Therefore
\[
 \frac{1+h_{T,\Delta}}{1+H_T}
 =1-\frac{k_T}{1+H_T}=1+o(1) \tag{8}
\]
uniformly over all deterministic \(\Delta_T\geq0\).  The bias-radius definition then gives
\[
 B_{T,\Delta}=O\!\left((1+H_T)^{-\beta}\right). \tag{9}
\]
Equations (2), (4), and (9) bound the concentration-branch radius by a constant multiple of
\(
r_{T,\Delta}=\max\{T^{-q},(1+H_T)^{-\beta}\}
\).
Equations (5)--(7) and (9) give the same bound for the Gaussian-branch radius.  Intersection with \([-2,2]\) cannot increase length.  On \(E_T^c\), the interval has length four, and \(\mathrm{lem:quota\text{-}success}\) gives
\[
 4P(E_T^c)\leq16T^{-4}=o(T^{-q})=o(r_{T,\Delta}). \tag{10}
\]
Taking expectations proves the stated uniform upper bound.

\emph{Sampling converse.}
Fix the same deterministic policy in both experiments:
\[
 p_t:=\pi_t(1)=ct^{-\alpha},
 \qquad \pi_t(0)=1-p_t. \tag{11}
\]
It is admissible for \(\mathfrak M_T(\alpha,\beta,L,c,C,\varepsilon)\): the rare-arm lower bound holds with equality, the upper bound follows from \(c\leq C\), and deterministic randomization with probability \(p_t\) satisfies sequential randomization.  Notice that no upper allocation bound is used after this membership check.  Set both delays to zero and set \(Y_t(0)=0\).  Under \(P_+\) and \(P_-\), respectively, let the independent and identically distributed arm-one outcome be Rademacher with means \(d_T\) and \(-d_T\), where
\[
 Q_T=\sum_{t=1}^T p_t,
 \qquad d_T=\lambda Q_T^{-1/2},
 \qquad 0<\lambda<\frac{1-2\eta}{4\sqrt2}. \tag{12}
\]
Because \(\alpha<1\), integral comparison gives \(Q_T\asymp T^{1-\alpha}\), and hence \(d_T\asymp T^{-q}\) and \(d_T<1\) for all sufficiently large \(T\).  The outcomes are bounded, the delays satisfy the polynomial envelope, and the two causal contrasts are
\(
\tau_+=d_T
\)
and
\(
\tau_-=-d_T
\).

The assignment and logged-propensity laws are common to the two experiments.  Given an arm-zero assignment, the observed-record laws agree; given an arm-one assignment, the Hellinger affinity of the two Rademacher laws is \(\sqrt{1-d_T^2}\).  Independence of the arrival vectors and the deterministic policy in (11) thus give the exact affinity of the observed experiments:
\[
 \begin{aligned}
 \operatorname{Aff}(P_+^O,P_-^O)
 &=\prod_{t=1}^T
   \left[1-p_t\{1-\sqrt{1-d_T^2}\}\right]\\
 &\geq1-\sum_{t=1}^T p_td_T^2
 =1-\lambda^2. 
 \end{aligned} \tag{13}
\]
Here \(1-\sqrt{1-x}\leq x\) for \(x\in[0,1]\), and
\(
\prod_t(1-x_t)\geq1-\sum_t x_t
\)
for \(x_t\in[0,1]\).  For any two dominated laws with densities \(p\) and \(q\), Cauchy--Schwarz gives
\[
 \begin{aligned}
 \|P-Q\|_{\mathrm{TV}}
 &=\frac12\int|\sqrt p-\sqrt q|(\sqrt p+\sqrt q)\\
 &\leq\frac12
 \left\{\int(\sqrt p-\sqrt q)^2
       \int(\sqrt p+\sqrt q)^2\right\}^{1/2}
 =\sqrt{1-\operatorname{Aff}(P,Q)^2}.
 \end{aligned} \tag{14}
\]
Equations (13)--(14) imply
\[
 \|P_+^O-P_-^O\|_{\mathrm{TV}}
 \leq\sqrt{1-(1-\lambda^2)^2}
 \leq\sqrt2\lambda. \tag{15}
\]

Let \(J_{T,\Delta}\) have the asserted uniform coverage; equivalently, allow a sequence \(\gamma_T\to0\) such that each of these two laws has coverage at least \(1-\eta-\gamma_T\).  Transferring the coverage of \(\tau_-\) from \(P_-^O\) to \(P_+^O\), then applying the union bound and (15), yields
\[
 P_+^O\{\tau_+,\tau_-\in J_{T,\Delta}\}
 \geq1-2\eta-2\gamma_T-\sqrt2\lambda. \tag{16}
\]
The right-hand side is eventually bounded below by a positive constant depending only on \(\eta\).  On this event the interval has length at least \(2d_T\).  Thus (12) and \(Q_T\asymp T^{1-\alpha}\) imply
\[
 \sup_{(P,\Pi,T)\in\mathfrak M_T}
 \mathbb E_{P,\Pi}|J_{T,\Delta}|
 \geq c_{\eta,1}T^{-q} \tag{17}
\]
for a constant \(c_{\eta,1}>0\).  This calculation uses a single fixed admissible policy; it neither requires nor exploits the upper allocation envelope beyond placing that policy in the displayed subclass.

\emph{Hidden-tail converse.}
Retain the policy (11) and put
\[
 \delta_T=\min\{1/4,L(1+H_T)^{-\beta}\}. \tag{18}
\]
For each recruitment time and arm, draw a Bernoulli indicator of probability \(\delta_T\), independently over recruitment times.  On failure set \((Y_t(a),D_t(a))=(0,0)\); on success set \(D_t(a)=H_T+1\).  Under \(P_+\), give a successful arm-one observation outcome \(+1\) and a successful arm-zero observation outcome \(-1\); under \(P_-\), reverse these signs.  Any dependence between the two arms within an arrival vector may be fixed identically under the two laws.  The arrival vectors are independent and identically distributed over \(t\), and the outcomes are bounded.  For every integer \(u\geq0\),
\[
 P_\pm\{D_t(a)>u\}
 =\begin{cases}
 \delta_T,&0\leq u\leq H_T,\\
 0,&u\geq H_T+1,
 \end{cases}
 \leq\min\{1,L(1+u)^{-\beta}\}. \tag{19}
\]
Thus the laws belong to the model class and all delays are finite.

No successful atom is revealed by terminal time \(H_T\).  The immediate zero records, nonrevelation indicators, assignments, and logged propensities therefore have exactly the same joint law under \(P_+\) and \(P_-\), so the complete observed laws of \(O_{T,\Delta}\) coincide.  Their eventual causal contrasts nevertheless are
\[
 \tau_+=2\delta_T,
 \qquad \tau_-=-2\delta_T. \tag{20}
\]
Under the common observed law, uniform coverage and a union bound give
\[
 P\{\tau_+,\tau_-\in J_{T,\Delta}\}
 \geq1-2\eta-2\gamma_T. \tag{21}
\]
On this event the interval length is at least \(4\delta_T\).  Since \(H_T\geq T\to\infty\), equation (18) equals \(L(1+H_T)^{-\beta}\) eventually.  Hence
\[
 \sup_{(P,\Pi,T)\in\mathfrak M_T}
 \mathbb E_{P,\Pi}|J_{T,\Delta}|
 \geq c_{\eta,2}(1+H_T)^{-\beta} \tag{22}
\]
for some \(c_{\eta,2}>0\).  The witness deliberately couples the nonzero outcome to the long delay, so this tail converse does not assert a lower bound within an outcome--delay-independence subclass.

The affinity identity (13) and the indistinguishability identity preceding (20) hold at every finite horizon, including \(T\leq4\); thus the same-experiment stress test does not condition on quota success, propensity convergence, or recruitment success.  Combining (17) and (22), and using
\[
 \max\{c_{\eta,1}x,c_{\eta,2}y\}
 \geq\min\{c_{\eta,1},c_{\eta,2}\}\max\{x,y\},
\]
proves the lower bound at
\(
r_{T,\Delta}=\max\{T^{-q},(1+H_T)^{-\beta}\}
\).
When \(\Delta_T=0\), \(H_T=T\), which gives
\(
\max\{T^{-(1-\alpha)/2},T^{-\beta}\}
\), as claimed.

%%% FIELD construction-one-sided-primary-class
\(\mathfrak M_T^{\mathrm{low}}(\alpha,\beta,L,c,\varepsilon)=\{(P,\Pi,T):\text{ass:iid-arrivals, ass:bounded-outcomes, ass:polynomial-delay-envelope, ass:sequential-randomization, ass:one-sided-forced-exploration, and ass:control-floor hold at index }T\}\).
